
\documentclass[a4paper,11pt]{article}
\usepackage[a4paper, margin=8em]{geometry}

% usa i pacchetti per la scrittura in italiano
\usepackage[french,italian]{babel}
\usepackage[T1]{fontenc}
\usepackage[utf8]{inputenc}
\frenchspacing 

% usa i pacchetti per la formattazione matematica
\usepackage{amsmath, amssymb, amsthm, amsfonts}

% usa altri pacchetti
\usepackage{gensymb}
\usepackage{hyperref}
\usepackage{standalone}

% imposta il titolo
\title{Appunti Comunicazioni Numeriche}
\author{Luca Seggiani}
\date{2026}

% imposta lo stile
% usa helvetica
\usepackage[scaled]{helvet}
% usa palatino
\usepackage{palatino}
% usa un font monospazio guardabile
\usepackage{lmodern}

\renewcommand{\rmdefault}{ppl}
\renewcommand{\sfdefault}{phv}
\renewcommand{\ttdefault}{lmtt}

% disponi il titolo
\makeatletter
\renewcommand{\maketitle} {
	\begin{center} 
		\begin{minipage}[t]{.8\textwidth}
			\textsf{\huge\bfseries \@title} 
		\end{minipage}%
		\begin{minipage}[t]{.2\textwidth}
			\raggedleft \vspace{-1.65em}
			\textsf{\small \@author} \vfill
			\textsf{\small \@date}
		\end{minipage}
		\par
	\end{center}

	\thispagestyle{empty}
	\pagestyle{fancy}
}
\makeatother

% disponi teoremi
\usepackage{tcolorbox}
\newtcolorbox[auto counter, number within=section]{theorem}[2][]{%
	colback=blue!10, 
	colframe=blue!40!black, 
	sharp corners=northwest,
	fonttitle=\sffamily\bfseries, 
	title=~\thetcbcounter: #2, 
	#1
}

% disponi definizioni
\newtcolorbox[auto counter, number within=section]{definition}[2][]{%
	colback=red!10,
	colframe=red!40!black,
	sharp corners=northwest,
	fonttitle=\sffamily\bfseries,
	title=~\thetcbcounter: #2,
	#1
}

% U.D.M
\newcommand{\amp}{\ensuremath{\mathrm{A}}}
\newcommand{\volt}{\ensuremath{\mathrm{V}}}
\newcommand{\meter}{\ensuremath{\mathrm{m}}}
\newcommand{\second}{\ensuremath{\mathrm{s}}}
\newcommand{\farad}{\ensuremath{\mathrm{F}}}
\newcommand{\henry}{\ensuremath{\mathrm{H}}}
\newcommand{\siemens}{\ensuremath{\mathrm{S}}}

% circuiti
\usepackage{circuitikz}
\usetikzlibrary{babel}

% disegni
\usepackage{pgfplots}
\pgfplotsset{width=10cm,compat=1.9}

% disponi codice
\usepackage{listings}
\usepackage[table]{xcolor}

\lstdefinestyle{codestyle}{
		backgroundcolor=\color{black!5}, 
		commentstyle=\color{codegreen},
		keywordstyle=\bfseries\color{magenta},
		numberstyle=\sffamily\tiny\color{black!60},
		stringstyle=\color{green!50!black},
		basicstyle=\ttfamily\footnotesize,
		breakatwhitespace=false,         
		breaklines=true,                 
		captionpos=b,                    
		keepspaces=true,                 
		numbers=left,                    
		numbersep=5pt,                  
		showspaces=false,                
		showstringspaces=false,
		showtabs=false,                  
		tabsize=2
}

\lstdefinestyle{shellstyle}{
		backgroundcolor=\color{black!5}, 
		basicstyle=\ttfamily\footnotesize\color{black}, 
		commentstyle=\color{black}, 
		keywordstyle=\color{black},
		numberstyle=\color{black!5},
		stringstyle=\color{black}, 
		showspaces=false,
		showstringspaces=false, 
		showtabs=false, 
		tabsize=2, 
		numbers=none, 
		breaklines=true
}

\lstdefinelanguage{javascript}{
	keywords={typeof, new, true, false, catch, function, return, null, catch, switch, var, if, in, while, do, else, case, break},
	keywordstyle=\color{blue}\bfseries,
	ndkeywords={class, export, boolean, throw, implements, import, this},
	ndkeywordstyle=\color{darkgray}\bfseries,
	identifierstyle=\color{black},
	sensitive=false,
	comment=[l]{//},
	morecomment=[s]{/*}{*/},
	commentstyle=\color{purple}\ttfamily,
	stringstyle=\color{red}\ttfamily,
	morestring=[b]',
	morestring=[b]"
}

% disponi sezioni
\usepackage{titlesec}

\titleformat{\section}
	{\sffamily\Large\bfseries} 
	{\thesection}{1em}{} 
\titleformat{\subsection}
	{\sffamily\large\bfseries}   
	{\thesubsection}{1em}{} 
\titleformat{\subsubsection}
	{\sffamily\normalsize\bfseries} 
	{\thesubsubsection}{1em}{}

% disponi alberi
\usepackage{forest}

\forestset{
	rectstyle/.style={
		for tree={rectangle,draw,font=\large\sffamily}
	},
	roundstyle/.style={
		for tree={circle,draw,font=\large}
	}
}

% disponi algoritmi
\usepackage{algorithm}
\usepackage{algorithmic}
\makeatletter
\renewcommand{\ALG@name}{Algoritmo}
\makeatother

% disponi numeri di pagina
\usepackage{fancyhdr}
\fancyhf{} 
\fancyfoot[L]{\sffamily{\thepage}}

\makeatletter
\fancyhead[L]{\raisebox{1ex}[0pt][0pt]{\sffamily{\@title \ \@date}}} 
\fancyhead[R]{\raisebox{1ex}[0pt][0pt]{\sffamily{\@author}}}
\makeatother

\begin{document}
% sezione (data)
\section{Lezione del 09-04-26}

% stili pagina
\thispagestyle{empty}
\pagestyle{fancy}

% testo
Abbiamo introdotto nella scorsa lezione le \textit{trasformazioni} di variabile aleatoria, ed un primo strumento per ricavare la distribuzione di probabilità di variabili trasformate, ovvero il metodo della \textit{funzione di distribuzione} (12.2.1). Vediamo adesso un risultato più specifico, ma che rende questo calcolo molto più veloce.

\subsubsection{Trasformazioni di variabili continue}
Diamo un importante teorema riguardante la trasformazione di variabili continue:
\begin{theorem}{Teorema fondamentale delle trasformazioni di variabili continue}
	Se le variabili aleatorie $X$ e la sua trasformata $Y = g(X)$ sono entrambe \textit{continue}, è possibile esprimere direttamente la distribuzione di probabilità di $Y$ mediante quella di $X$, come:
$$
f_Y (y) = \frac{f_X (x)}{g'(x)} \Bigg|_{x = g^{-1}(y)}
$$
\end{theorem}
Se la funzione fosse monotona decrescente a denominatore ci sarebbe il valore assoluto di $g'(x)$ (una d.d.p. è sempre non negativa).

Questo risultato deriva dal fatto che, posto $y = g(x)$:
$$
f_Y (y) dy = f_X (x) dx, \quad
\begin{cases}
f_Y (y) dy = P\left( \{ y < Y \leq y + dy \} \right) \\
f_X (x) dx = P\left( \{ x < X \leq x + dx \} \right) \\
\end{cases}
$$
in quanto per gli eventi vale che:
$$
\{ y < Y \leq y + dy \} \equiv \{ x < X \leq x + dx \} \implies
P\left( \{ y < Y \leq y + dy \} \right) = P\left( \{ x < X \leq x + dx \} \right)
$$
cioè sono uguali, ovvero composti dagli stessi risultati (chiaramente trasformati da $X$ in $Y$, attraverso $g$).
Riarrangiando, si ottiene:
$$
f_Y (y) = f_X (x) \frac{dx}{dy} = \frac{f_X (x)}{g'(x)} \Bigg|_{x = g^{-1}(y)}
$$
dove chiaramente si deve imporre $x = g^{-1}(y)$ da $y = g(x)$ posto prima. \hfill $\square$

Questo teorema vale anche per intervalli disgiunti (nel caso di funzioni continue non iniettive), in quanto se in generale $g(X)$ è una funzione continua, che non è costante in alcun intervallo (altrimenti anche se $X$ fosse continua non lo sarebbe $Y$), si ha:
$$
\{ y < Y \leq y + dy \} = \sum_{i = 1}^N \{ x_i < X < x_i + d x_i \}
$$
su $N$ intervalli, per cui:
$$
f_Y (y) dy = \sum_{i = 1}^N f_X (x_i) dx_1
$$
con gli opportuni valori assoluti per punti decrescenti di $f_X (x)$. Applicando il teorema fondamentale si ricava quindi:
$$
f_Y (y) dy = \sum_{i = 1}^N \frac{f_X (x_i)}{g'(x_i)} \Bigg|_{x = g^{-1}(y)} 
$$
notiamo quindi che si ha per le controimmagini degli $y$:
$$
\{x_i : g(x_i) = y\}
$$
cioè gli $x_i$ sono le soluzioni dell'equazione:
$$
y = g(x)
$$
con $y$ fisso (nel punto in cui vogliamo calcolare $f_Y (y)$). Chiaramente il numero $N$ di $x_i$ varia al variare di $y$ (per una varietà di funzioni non inettive). Inoltre, se per un certo $y^*$ l’equazione $y = g(x)$ non ha soluzione, risulterà:
$$
f_Y (y^*) = 0
$$

\subsection{Indici caratteristici di una distribuzione}
La funzione di distribuzione $F_X(x)$ oppure la densità di probabilità $f_X(x)$ o la massa di probabilità $p_X(x)$ forniscono una descrizione statistica completa della variabile aleatoria $X$, cioè il massimo di informazione che si può avere sul comportamento statistico dei valori assunti da $X$. 

In molti casi, però, non è possibile arrivare a una conoscenza così completa riguardo a un problema aleatorio che si sta trattando. In questo caso, se non si riesce ricavare la legge di distribuzione, allora ci accontentiamo di alcuni indici caratteristici (parametri statistici semplificati) che forniscono informazioni parziali sulla legge di distribuzione. Questi sono:
\begin{itemize}
	\item Il \textbf{valor medio} (o \textit{valore atteso});
	\item La \textbf{varianza};
	\item La \textbf{deviazione standard};
	\item Il \textbf{valore quadratico medio}.
\end{itemize}

\subsubsection{Valore medio di una variabile aleatoria discreta}
Per una variabile aleatoria discreta, il valore medio è definito come:
$$
\eta_X = E\left[X\right] = \sum_i x_i \cdot p_X (x_i)
$$
La lettera $E$ viene dall'inglese \textit{Expectation}, cioè aspettativa. Il valor medio di una variabile aleatoria $X$ è compreso tra il più piccolo ed il più grande valore assunto da $X$, e fornisce un’indicazione sulla posizione attorno alla quale si addensano i valori della $X$ (\textit{indice di posizione}). Notiamo che non è detto che il valor medio $\eta_X$ di una variabile aleatoria discreta $X$ corrisponda con alcun valore assunto da $X$.

\subsubsection{Varianza di una variabile aleatoria discreta}
La varianza (e vedremo a breve la deviazione standard) misurano lo scostamento dei valori di una variabile aleatoria $X$ dal suo valor medio, ovvero sono entrambi un \textit{indice di dispersione} della variabile aleatoria.

In particolare la varianze è definita come:
$$
\sigma^2_X = E \left[ (X - \eta_X)^2 \right] = \sum_i (x_i - \eta_X)^2 p_X (x_i)
$$

\subsubsection{Deviazione standard di una variabile aleatoria discreta}
La deviazione standard è quindi definita direttamente in funzione della varianza:
$$
\sigma_X = \sqrt{\sigma^2_X} = \sqrt{ E \left[ (X - \eta_X)^2 \right] } = \sqrt{ \sum_i (x_i - \eta_X)^2 p_X (x_i) }
$$
La deviazione standard è quindi anch'essa una misura dello scostamento dei valori di $X$ dal valor medio.

\subsubsection{Valore quadratico medio di una variabile aleatoria discreta}
Vediamo infine il valore quadratico medio di una variabile discreta $X$, definito come la sua potenza media statistica:
$$
P_X = E\left[ X^2 \right] = \sum_i x_i^2 p_X(x_i)
$$
Il valore quadratico medio e la varianza sono legati da una semplice relazione:
$$
\sigma^2_X = E \left[ (X - \eta_X)^2 \right] = E \left[ X^2 - 2\eta_X + \eta_X^2 \right] = E \left[ X^2 \right] - 2 \eta_X E \left[ X \right] + \eta_X^2 = P_X - \eta_X^2
$$
dove i passaggi sono giustificati dalla proprietà di \textit{linearità} dell'aspettazione. 

\subsubsection{Valore medio di una variabile aleatoria continue}
Per una variabile aleatoria continua, il valore medio è definito come:
$$
\eta_X = E\left[X\right] = \int_{-\infty}^{\infty} x f_X(x) \, dx
$$
cioè l'analogo continuo del valor medio di variabili discrete.

Si nota immediatamente che, nel caso in cui la funzione densità di probabilità $f_X(x)$ di una certa variabile aleatoria continua è pari, il valor medio è immediatamente nullo. Infatti, in:
$$
\int_{-\infty}^{\infty} x f_X(x) \, dx = 0
$$
si ha $x$ dispari, $f_X(x)$ pari, e quindi i rami a sinistra e a destra di $x_0 = 0$ che si annullano.

Questo si generalizza a simmetrie rispetto a valori $a$, per cui il valor medio vale $a$:
$$
f_X(a - x) = f_X(a + x) \implies \eta_X = a
$$
in quanto:
$$
\int_{-\infty}^{\infty} (a - x) f_X(x) \, dx = 0 = a \int_{-\infty}^\infty f_X(x) \, dx - \int_{-\infty}^{\infty} x f_X(x) \, dx = a - \eta_x
$$
da cui la tesi. \hfill $\square$.

\subsubsection{Varianza di una variabile aleatoria continua}
La varianza per una variabile aleatoria continua è definita come:
$$
\sigma^2_X = E \left[ (X - \eta_X)^2 \right] = \int_{-\infty}^{\infty} (x - \eta_X)^2 f_X (x) \, dx 
$$

\subsubsection{Deviazione standard di una variabile aleatoria cotinua}
La deviazione standard è quindi definita direttamente in funzione della varianza anche per le variabili aleatorie continue:
$$
\sigma_X = \sqrt{\sigma^2_X} = \sqrt{ E \left[ (X - \eta_X)^2 \right] } = \sqrt{ \int_{-\infty}^{\infty} (x - \eta_X)^2 f_X (x) \, dx }
$$

\subsubsection{Valore quadratico medio di una variabile aleatoria continua}
Vediamo infine il valore quadratico medio anche per una variabile continua $X$, definito come la sua potenza media statistica:
$$
P_X = E\left[ X^2 \right] = \int_{-\infty}^{\infty} x^2 f_X(x) \, dx
$$

\subsection{Teorema dell'aspettazione}
Se $Y = g(X)$ è possibile ricavare il valore medio di $Y$ direttamente dalla distribuzione di $X$.

Consideriamo il caso in cui $Y$ ed $X$ sono variabili discrete. In tal caso vale:
$$
\eta_Y = E(Y) = \sum_k y_k P(Y = y_k) = y_1 P \{ Y = y_1 \} + y_2 P \{ Y = y_2 \} + ...
$$
Il primo termine di questa sommatoria si potrà scrivere come:
$$
y_1 P \{ Y = y_1 \} = y_1 \sum_{G(y_1)} P \{ X = x_i \} = \sum_{G(y_1)} g(x_i) P(X = x_i)
$$
dove la sommatoria è estesa a tutti gli $x_i$ per i quali $g(x_i) = y_i$, ovvero:
$$
G(y_1) = \{ x_i : g(x_i) = y_i \}
$$
Notiamo che gli $G(y_1), ... G(y_n)$ formano una partizione di $\mathbb{R}$, per cui $E\left[Y\right]$ può essere espressa come:
$$
E\left[Y\right] = E\left[g(X)\right] = \sum_i g(x_i) P(X = x_i)
$$

Da questo si ha il teorema:
\begin{theorem}{Teorema dell'aspettazione per variabili aleatore discrete}
L'aspettazione di una variabile aleatoria $Y$ trasformata da una variabile aleatoria $X$ come $Y = g(X)$ si può calcolare come:
$$
E\left[Y\right] = E\left[g(X)\right] = \sum_i g(x_i) P(X = x_i)
$$
\end{theorem}

Notiamo che la definizione di varianza $X$ è un caso particolare di questo teorema:
$$
\sigma^2_X = E \left[ (X - \eta_X)^2 \right] = \sum_i (x_i - \eta_X)^2 p_X (x_i)
$$

Infine, diamo lo stesso teorema nel caso continuo (la dimostrazione sarà analoga): 
\begin{theorem}{Teorema dell'aspettazione per variabili aleatore continue}
L'aspettazione di una variabile aleatoria $Y$ trasformata da una variabile aleatoria $X$ come $Y = g(X)$ si può calcolare come:
$$
E\left[Y\right] = E\left[g(X)\right] = \int_{-\infty}^{\infty} g(x) f_X(x) \, dx 
$$
\end{theorem}

\subsubsection{Proprietà dell'operatore aspettazione}
Abbiamo anticipato alcune proprietà dell'operatore aspettazione, fra cui la \textit{linearità}. Vediamole tutte in ordine:
\begin{enumerate}
\item \begin{theorem}{Linearità dell'aspettazione}
Vale che l'operatore di aspettazione è lineare:
$$
\begin{cases}
	E\left[ g(X) + c \right] = E\left[ g(X) \right] + c \\
	E\left[ cg(X) \right] = c E\left[ g(X) \right] \\
	E\left[ g(X) + h(X) \right] = E\left[ g(X) \right] E\left[ h(X) \right]
\end{cases}
$$
\end{theorem}
Questo si dimostra direttamente dal teorema dell'aspettazione nel caso delle variabili discrete:
$$
E\left[ g(X) + c \right] = \sum_i \left( g(x_i) + c \right) P(X = x_i) = \sum_ig(x_i) P(X = x_i) + c P(X = x_i) 
$$
$$
= \sum_ig(x_i) P(X = x_i) + c = E\left[ g(X) \right] + c
$$
per la prima proprietà, quindi:
$$
E\left[cg(X)\right] = c \sum_i g(x_i) P(X = x_i) = c E \left[ g(X) \right]
$$
e infine:
$$
E \left[ g(X) + h(X) \right] = \sum_i \left[ g(x_i) + h(x_i) \right] P (X = x_i)
$$
$$
= \sum_i g(x_i) P (X = x_i) + \sum_i h(x_i) P (X = x_i) = E\left[ g(X) \right] E\left[ h(X) \right]
$$
da cui la tesi. \hfill $\square$

Il caso più generale sarà quindi:
$$
E \left[ \sum_i \alpha_i g_i (X) \right] = \sum_i \alpha_i E \left[ g_i(X) \right]
$$
\end{enumerate}

\end{document}
